\section{From Ward identities to invariant curvature coefficients}
\label{sec:multileg-curvature-lift}
The leading Taylor terms of a gauge-invariant action satisfy a
Ward identity in every field argument. We now turn these identities
into an exponentially summable expansion in curvature variables.
The momenta of the arguments are constrained to sum to zero, so
one cannot simply apply independent one-variable divisions to all
of them. The following construction resolves that dependence before
choosing the invariant, real coefficients of the path expansion.

\subsection{A single exterior division with auxiliary constraints}
Put $q_a(p)=e^{ip_a}-1$. On alternating coordinate tensors in four
dimensions let $E_a$ denote exterior multiplication by the $a$th
basis vector and let $\iota_a=E_a^{\mathsf T}$. Write
\[
 \delta(p)=\sum_{a=1}^4 q_a(p)\iota_a.
\]
Thus $\delta$ lowers exterior degree by one and $\delta^2=0$.
Tensor factors other than the one on which this operator acts
are coefficient spaces. Operators on distinct factors commute.
All transposes here are algebraic transposes, without complex
conjugation. This convention is appropriate for the multilinear
Taylor coefficients of a scalar action.

A periodic coefficient on $d$ complex momentum coordinates has
strip bound $(\sigma,A)$ if it is holomorphic on
$|\operatorname{Im}p_j|<\sigma$ and has absolute value at most $A$
there. Tensor bounds mean this entry by entry. Say that it has
order $s$ at zero if its restriction to each real ray $p=tu$ is
$O(t^s)$ as $t\to0$. For an analytic function this is equivalent
to the vanishing of all homogeneous Taylor polynomials of degree
less than $s$: restrict the Taylor expansion to each ray and use
that a polynomial vanishing on every real vector is zero.

The division from Section~\ref{sec:curvature-horizon-producer}
has the following strip version. For each coordinate $j$, there
is a periodic holomorphic $D_j f$ such that
\begin{equation}\label{eq:multileg-division}
 (e^{ip_j}-1)D_jf(p)=f(p)-f(p|_{p_j=0}),
 \qquad
 |D_jf|_{\sigma/4}\le A\,d_\sigma,
\end{equation}
where one may take
\[
 d_\sigma=\frac4\sigma
       \left(\frac{2(1+\sigma/4)}{\sigma/4}\right)^d.
\]
Indeed, Fourier expansion gives coefficients $\widehat f(z)$ with
\[
 \sum_z e^{(\sigma/2)|z|_1}|\widehat f(z)|
 \le A\left(\frac{2(1+\sigma/4)}{\sigma/4}\right)^d.
\]
Replace each monomial in coordinate $j$ by the finite geometric
sum displayed in that earlier section. It has at most $|z|_1$
terms, all with displacement norm at most $|z|_1$. On the smaller
strip its absolute value is at most
$|z|_1e^{(\sigma/4)|z|_1}$, which is bounded by
$(4/\sigma)e^{(\sigma/2)|z|_1}$. This proves normal convergence,
the bound, and the identity, including the zeros of the divisor.
If $f$ has order $s$, this division has order at least
$\max(s-1,0)$. To verify this, the right side of
\eqref{eq:multileg-division} has order $s$, while the lowest
nonzero homogeneous term of the divisor is $ip_j$. A nonzero
lowest homogeneous term of the quotient therefore gains exactly
one degree upon multiplication.

Suppose a tensor $f$ of exterior degree $r$ is a cycle,
$\delta f=0$. Set $f_0=f|_{p_a=0}$ and
\[
 B_a=E_aD_af,\qquad R_a=(I-E_a\iota_a)f_0.
\]
The Cartan identity $\delta E_a+E_a\delta=q_a I$ gives
\begin{equation}\label{eq:multileg-axis}
                  f=\delta B_a+R_a.
\end{equation}
For example, where $q_a\ne0$ the cycle identities for $f$ and
$f_0$ imply $\delta D_af=-\iota_af_0$; substitution gives the
displayed formula. Continuity extends it to $q_a=0$.
The residual is still a cycle, has zero entries on faces containing
$a$, and evaluates all other entries at $p_a=0$.
Successive residual operations in the four coordinates therefore
leave zero when $r>0$. When $r=0$, they leave the value at $p=0$,
so the same conclusion holds if that value vanishes. Adding the
four lifts proves the following fact.

\begin{lemma}\label{lem:multileg-single}
Let $r,s\in\mathbb N$. Let $f$ have strip bound $(\sigma,A)$,
order $s$, and exterior degree $r$. Assume $\delta f=0$ and either $r>0$ or $f|_{p=0}=0$.
Then $f=\delta B$, where $B$ has degree $r+1$, strip bound
$(\sigma/4,4Ad_\sigma)$, and order at least $\max(s-1,0)$.
Here other momentum coordinates may remain as parameters, and
$p=0$ in the degree-zero hypothesis means only that the four
coordinates of the divided momentum vanish.

If, in addition, for each exterior component the vector of
coefficients in the other tensor factors belongs to a fixed
linear subspace of that other-factor coefficient space, depending
only on the other momenta, the same membership holds separately
for every exterior component of $B$ at every real momentum.
\end{lemma}
\begin{proof}
Only the last assertion remains. Away from $q_a=0$, the quotient
in \eqref{eq:multileg-division} is a scalar multiple of the
difference of two vectors in that subspace. On $q_a=0$, vary just
$p_a$ and take the limit. The subspace is fixed during this limit
and is closed, since the coefficient space is finite dimensional.
Exterior multiplication, subtraction, and coordinate restriction
preserve membership as well. Thus each of the four lifts and
their sum preserve the constraint. Each residual retains the
original strip bound and order; summing four lifts accounts for
the factor four, with only one loss of Taylor order.
\end{proof}

\subsection{Restoring the dependent momentum constraint}
There are $t$ independent momenta $p_1,\ldots,p_t\in\mathbb C^4$,
and we set $p_0=-\sum_{j=1}^t p_j$. Write $\delta_j=\delta(p_j)$
acting on tensor factor $j$. A Ward identity on that factor is
$\delta_j f=0$. We first allow the zeroth factor to have degree
$r+1$ and all independent factors to have degree one.

One can solve the zeroth-factor cycle while preserving all Ward
identities except, provisionally, the first. Exchange $p_1$ and
$p_0$ by the linear involution
\[
 (p_1,p_2,\ldots,p_t)\longmapsto
             (-p_1-\cdots-p_t,p_2,\ldots,p_t).
\]
It preserves the integer periodicity lattice and sends a sufficiently
narrow strip into the original strip; width $\sigma/(t+1)$ is
sufficient for $t\ge1$. Apply Lemma~\ref{lem:multileg-single}
in the new first variable and exchange back. The Ward subspaces
for factors $2,\ldots,t$ do not depend on this variable, so they
are retained. The cost is a fixed multiple of the old cost,
the width remains strictly positive, and Taylor order drops by
at most one. Denote this operation simply as a zeroth-factor lift.
Likewise, if a coefficient has no first tensor factor and vanishes
when $p_1=0$, the degree-zero case of the lemma adds a degree-one
first factor whose contraction is the original coefficient,
while retaining the remaining Ward identities.

\begin{lemma}[Dimension shift]\label{lem:multileg-shift}
Fix $t,r\ge0$ with $r+1+2t>4$ and $\sigma>0$.
There are $\eta>0$ and $K<\infty$, depending only on these data,
such that the following holds for every $A\ge0$.
If $u$ has strip bound $(\sigma,A)$, degrees
$(r+1,1,\ldots,1)$, and $\delta_j u=0$ for $j=0,\ldots,t$,
there is $v$ of degrees $(r+2,1,\ldots,1)$ with strip bound
$(\eta,KA)$ satisfying
\[
                 \delta_0v=u,\qquad \delta_jv=0\quad(1\le j\le t).
\]
\end{lemma}
\begin{proof}
Induct on $t$, allowing $r$ to vary. At $t=0$ the degree $r+1$
exceeds four, so $u=0$ and one takes $v=0$.
For $t\ge1$, first lift the zeroth cycle, temporarily forgetting
the first Ward identity. This gives $X$ with $\delta_0X=u$
and all remaining Ward identities. Contract its first factor:
\[
                  V=\delta_1X,\qquad \delta_0V=\delta_1u=0.
\]
A second zeroth-factor lift gives $Y$ of zeroth degree $r+3$
with $\delta_0Y=V$ and the remaining Ward identities.
Restrict $Y$ to $p_1=0$, obtaining $Z$. Since $q(0)=0$,
$V|_{p_1=0}=0$, so $Z$ is a zeroth-factor cycle in $t-1$
independent variables. Its zeroth degree is $r+3$, and
\[
                  r+3+2(t-1)=r+1+2t>4.
\]
The induction hypothesis gives $W$ of degree $r+4$ with
$\delta_0W=Z$, retaining every remaining Ward identity.
Extend $W$ constantly in $p_1$, writing the extension as $W'$,
and set
\[
                         L=Y-\delta_0W'.
\]
Here $\delta_0$ uses the full momentum $-p_1-\cdots-p_t$.
Thus $\delta_0L=V$ by $\delta_0^2=0$, while $L|_{p_1=0}=0$
by the induction identity. The degree-zero first-factor lift
now gives $F$ with $\delta_1F=L$ and all remaining Ward
identities. Finally put $v=X-\delta_0F$. Commutation on distinct
factors gives
\[
 \delta_0v=u,\qquad
 \delta_1v=V-\delta_0L=0.
\]
The other identities hold term by term. All operations are finite
sums, fixed linear changes of variables, restrictions, exterior
contractions, and the lifts already bounded above. At each stage
choose a smaller positive common width. Their number is determined
by $t,r$, and every bound remains linear in $A$. This gives $\eta$
and $K$ before the particular input tensor is chosen.
\end{proof}

Take now a tensor $f$ with $t+1\ge3$ degree-one factors and all
Ward identities. The dimension-shift lemma with $r=0$ lifts
factor zero while preserving the independent identities. Lift
each independent factor in turn by Lemma~\ref{lem:multileg-single}.
At each step the identities in the remaining factors are auxiliary
subspaces independent of the divided momentum; the already lifted
indices are just finite coefficient indices. We obtain $g$ with
all factors of degree two such that
\begin{equation}\label{eq:multileg-all-curls}
 f(p_1,\ldots,p_t)=\delta_0\delta_1\cdots\delta_t
                           g(p_1,\ldots,p_t).
\end{equation}
Its strip width and cost multiplier depend only on the original
width and the number of factors. In particular the cubic and
quartic applications have constants independent of the torus.

\subsection{The quadratic case and retained moments}
For two factors the dimension inequality fails. Suppose instead
that $f(p)$ has both Ward identities at momenta $(-p,p)$ and
order $s+2$ at zero. The preceding construction can be shortened
in a way that retains precisely $s$ orders after both lifts.
Lift factor zero to obtain $X$ of degrees $(2,1)$ and order
$s+1$, with $\delta_0X=f$. The contraction
$V=\delta_1X$ has order $s+2$, since the contraction symbol
vanishes to first order. It is a zeroth-factor cycle.
Lift it to obtain $Y$ of zeroth degree three and order $s+1$
with $\delta_0Y=V$.

There are no other momenta now. Hence $Y|_{p=0}=0$, and the
degree-zero first-factor lift gives $D$ of degrees $(3,1)$,
order $s$, and $\delta_1D=Y$. Set $Z=X-\delta_0D$.
Then
\[
       \delta_0Z=f,\qquad \delta_1Z=0,
                   \qquad Z\text{ has order }s+1.
\]
One last first-factor lift gives $g$ of degrees $(2,2)$ and
order $s$, with $\delta_0\delta_1g=f$. All bounds remain linear
in the initial strip cost, with a positive width chosen before
$f$ and $s$. For the loop remainder we use $s=2$: fourth-order
vanishing of the potential coefficient yields second-order
vanishing of its curvature coefficient.

\subsection{Real coefficients and invariance at each placement}
Fix a real orthonormal basis of the finite-dimensional Lie algebra.
The tensors now also have one colour index for each field factor.
Assume the original tensor satisfies
\[
             \overline{f(-p)}=f(p)
                 \quad\text{for real independent momenta }p,
\]
and that its real and imaginary colour tensors are invariant
under the simultaneous adjoint action of the gauge group.
The exterior lifts need not respect these properties, so we
restore them explicitly.

For each entry of $g$, take the real part of its Fourier
coefficients. The resulting real coefficient array has symbol
\[
                   g_{\mathbb R}(p)
                      =\tfrac12\bigl(g(p)+\overline{g(-p)}\bigr).
\]
Since $\overline{q(-p)}=q(p)$ for real $p$, contraction of
$g_{\mathbb R}$ is still $f$. If $g$ has order $s$, so does
$g_{\mathbb R}$. Fourier coefficient estimates on a smaller strip
give an exponentially weighted summation bound. More precisely,
from strip width $w$ and entry bound $B$ one obtains width $w/2$
and bound
\[
 B\left(\frac{2(1+w/4)}{w/4}\right)^{4t}.
\]

Let $\mathcal I$ be the subspace of invariant real colour tensors,
and choose a fixed linear projection $P$ onto $\mathcal I$.
Apply $P$ separately at each tuple of face directions and each
tuple of lattice displacements. This gives real coefficients
$R_{\mathbf a,\mathbf c}(z)$ that are invariant at every such
placement. Exterior contractions do not act on colour indices,
so they commute with $P$. Since $P$ fixes the real and imaginary
parts of $f$, the reconstructed potential coefficient remains
$f$. The weighted bound is multiplied by at most
\[
                         \sum_{i,j}|P_{ij}|<\infty.
\]
This is a fixed finite-dimensional projection; no assertion that
its norm is one is needed. It also preserves each retained order
of vanishing. Thus, for some $\eta>0$ and $C<\infty$ fixed by
width, arity, and gauge group, every face and colour entry obeys
\begin{equation}\label{eq:multileg-real-budget}
             \sum_{z\in\mathbb Z^{4t}}
                  e^{\eta|z|_1}|R_{\mathbf a,\mathbf c}(z)|\le CA.
\end{equation}

For two factors with retained order two, these coefficients have
both zero total mass and zero first moments, entry by entry:
\begin{equation}\label{eq:multileg-moments}
             \sum_z R(z)=0,\qquad \sum_z z_aR(z)=0
                                      \quad(a=1,\ldots,4).
\end{equation}
Indeed the symbol at zero is the first sum. Exponential summability
implies $\sum_z |z|_1|R(z)|<\infty$, so its derivative along the
$a$th real momentum coordinate is $i\sum_z z_aR(z)$. The symbol's
value and first derivatives vanish by its retained second order.
This proves both identities with absolute convergence.

\subsection{The invariant path expansion}
For each tuple $(\mathbf a,z)$ of oriented coordinate plaquettes
and displacements, put the zeroth plaquette at the origin and
plaquette $j$ at $z_j\in\mathbb Z^4$ for $1\le j\le t$.
Use the coefficient colour tensor $R_{\mathbf a,\cdot}(z)$ for
this atom. It is invariant by the preceding construction.
In the path convention of Section~\ref{sec:path-attachment-producer},
the atom's length is exactly
\[
                           |z|_1+4(t+1).
\]
The four edges of every plaquette, including the anchored one,
contribute to this length. Summing the weighted absolute
coefficients over the finite face and colour index sets and over
all displacements gives total atom budget at most
\[
 CA\,\#\{\text{face tuples}\}\,
       \#\{\text{colour tuples}\}\,e^{4\eta(t+1)}.
\]
An enumeration of this countable family changes neither its sum
nor its budget. In the two-factor case, filtering the enumeration
by a fixed pair of face directions and a fixed colour tuple
turns \eqref{eq:multileg-moments} into the required zero total
loop coefficient and zero first moment of the relative displacement
of the two legs. The latter displacement is $z_1-0$.

The potential coefficient reconstructed from these atoms is
exactly \eqref{eq:multileg-all-curls}, with the two-factor version
proved separately above. Absolute convergence permits
periodization and evaluation on every finite torus. This result
constructs the invariant coefficients from the specified analytic
Ward data. The subsequent stationary action calculation must
still supply those data, their vanishing orders, and the scalar
normalization of the recorded Taylor terms.
